% =============================================================================
% Präambel des Scientific Writing Kit
% =============================================================================
% Herstellerzone: Wird per /update erneuert. Eigene Anpassungen gehören nach
% latex/eigene-praeambel.tex (wird am Ende geladen, falls vorhanden).
%
% Engine: LuaLaTeX (MiKTeX, TeX Live, MacTeX). XeLaTeX (Tectonic) als Fallback.
% Alle \kit...-Makros kommen aus kit-daten.tex, das kit/werkzeuge/pdf.mjs aus
% arbeit/projekt.json erzeugt. Nie von Hand ändern, sondern die Config.
% =============================================================================

\usepackage{iftex}
\ifPDFTeX
  \errmessage{Bitte mit LuaLaTeX oder XeLaTeX kompilieren (node kit/werkzeuge/pdf.mjs)}
\fi

% --- Schrift ---------------------------------------------------------------
\usepackage{fontspec}
\ifx\kitSchrift\empty
  % Standard: Latin Modern, in jeder TeX-Distribution enthalten.
\else
  \setmainfont{\kitSchrift}
\fi

% --- Sprache ---------------------------------------------------------------
\ifkitEnglisch
  \usepackage[ngerman,main=english]{babel}
\else
  \usepackage[english,main=ngerman]{babel}
\fi
\usepackage[autostyle=true]{csquotes}

% --- Seite und Absatz ------------------------------------------------------
\usepackage[
  top=\kitRandOben, bottom=\kitRandUnten,
  inner=\kitRandInnen, outer=\kitRandAussen,
  includefoot, footskip=1.2cm
]{geometry}
\usepackage{setspace}
\setstretch{\kitZeilenabstand}
\usepackage{microtype}

% --- Kopf- und Fußzeile (KOMA) ---------------------------------------------
\usepackage[automark,headsepline=false]{scrlayer-scrpage}
\clearpairofpagestyles
\cfoot*{\pagemark}
\ohead{\headmark}
\setkomafont{pagehead}{\normalfont\small\itshape}
\setkomafont{pagenumber}{\normalfont}

% --- Überschriften ---------------------------------------------------------
% Serifen statt KOMA-Serifenlos, Größen wie in der erprobten Bachelor-Vorlage.
\setkomafont{disposition}{\normalfont\bfseries}
\ifdefined\chapter\setkomafont{chapter}{\Large}\fi
\setkomafont{section}{\large}
\setkomafont{subsection}{\normalsize}
\setkomafont{subsubsection}{\normalsize}
\ifdefined\chapter
  \RedeclareSectionCommand[beforeskip=0pt, afterskip=18pt]{chapter}
\fi
\RedeclareSectionCommand[beforeskip=-18pt plus -4pt minus -2pt, afterskip=6pt]{section}
\RedeclareSectionCommand[beforeskip=-12pt plus -3pt minus -2pt, afterskip=6pt]{subsection}
\setcounter{secnumdepth}{3}
\setcounter{tocdepth}{2}

% --- Verzeichnisse ---------------------------------------------------------
% Abstände im Inhaltsverzeichnis und Schutz gegen Überlappung langer Titel mit
% der Seitenzahl (erprobt in der Vorgängerarbeit).
\makeatletter
\renewcommand{\@pnumwidth}{2.2em}
\renewcommand{\@tocrmarg}{3.5em}
\makeatother
\ifdefined\chapter
  \DeclareTOCStyleEntry[beforeskip=10pt plus 1pt]{chapter}{chapter}
\fi
\DeclareTOCStyleEntry[beforeskip=2pt]{section}{section}

% --- Mathematik, Einheiten, Chemie -----------------------------------------
\usepackage{amsmath}
\usepackage{amssymb}
\usepackage{siunitx}
\ifkitEnglisch
  \sisetup{locale=UK, per-mode=symbol, separate-uncertainty=true}
\else
  \sisetup{locale=DE, per-mode=symbol, separate-uncertainty=true}
\fi
\usepackage[version=4]{mhchem}
\usepackage{chemfig}

% --- Abbildungen und Tabellen ----------------------------------------------
\usepackage{graphicx}
% Bilder liegen im Build-Ordner unter abbildungen/, Fallback: Projektwurzel.
\graphicspath{{./}{abbildungen/}{../../}}
\usepackage{xcolor}
\usepackage{booktabs}
\usepackage{longtable}
\usepackage{array}
\usepackage{tabularx}
\usepackage{calc}
\usepackage[font=small, labelfont=bf, format=hang]{caption}
\usepackage{float}
\usepackage{enumitem}
% Abbildungen und Tabellen durchgehend nummerieren (1, 2, 3 statt 3.1, 3.2).
\ifdefined\chapter
  \counterwithout{figure}{chapter}
  \counterwithout{table}{chapter}
  \counterwithout{equation}{chapter}
\fi
% Abstände zwischen Gleitobjekten und Text.
\setlength{\textfloatsep}{20pt plus 4pt minus 4pt}
\setlength{\intextsep}{20pt plus 4pt minus 4pt}
\setlength{\floatsep}{16pt plus 4pt minus 4pt}
% Nicht mehr als eine Seite nur mit Gleitobjekten erzwingen.
\renewcommand{\topfraction}{0.85}
\renewcommand{\textfraction}{0.1}
\renewcommand{\floatpagefraction}{0.75}

% --- Pandoc-Kompatibilität --------------------------------------------------
% Pandoc erzeugt diese Befehle in den Kapiteldateien.
\providecommand{\tightlist}{%
  \setlength{\itemsep}{0pt}\setlength{\parskip}{0pt}}
\makeatletter
\newsavebox\pandoc@box
\providecommand*\pandocbounded[1]{%
  \sbox\pandoc@box{#1}%
  \Gscale@div\@tempa{\textheight}{\dimexpr\ht\pandoc@box+\dp\pandoc@box\relax}%
  \Gscale@div\@tempb{\linewidth}{\wd\pandoc@box}%
  \ifdim\@tempb\p@<\@tempa\p@\let\@tempa\@tempb\fi
  \ifdim\@tempa\p@<\p@\scalebox{\@tempa}{\usebox\pandoc@box}%
  \else\usebox{\pandoc@box}%
  \fi}
\makeatother
\providecommand{\passthrough}[1]{#1}
\providecommand{\ul}[1]{\underline{#1}}
\providecommand{\st}[1]{\sout{#1}}
\usepackage[normalem]{ulem}
\usepackage{fancyvrb}
\usepackage{url}
\urlstyle{same}

% --- Literatur (biblatex + biber) ------------------------------------------
% Stil und Optionen setzt pdf.mjs aus arbeit/projekt.json -> zitation.stil.
\makeatletter
\edef\kit@biblatex{\noexpand\usepackage[backend=biber,\kitBibOptionen]{biblatex}}
\kit@biblatex
\makeatother
\kitBibNachladen
\addbibresource{literatur.bib}

% --- Entwurfs-Wasserzeichen ------------------------------------------------
\ifkitEntwurf
  \DeclareNewLayer[
    background, textarea,
    contents={%
      \raisebox{-0.55\layerheight}[0pt][0pt]{%
        \makebox[\layerwidth]{%
          \rotatebox{50}{\fontsize{80}{80}\selectfont\color{black!7}\kitWortEntwurf}}}}
  ]{kit.entwurf}
  \AddLayersToPageStyle{@everystyle@}{kit.entwurf}
\fi
% Platzhalter für noch fehlende Unterkapitel im Entwurfsmodus.
\newcommand{\kitPlatzhalter}[1]{%
  \section{#1}%
  \begin{center}
    \fcolorbox{black!30}{black!3}{\parbox{0.85\linewidth}{\centering\small\itshape
      \kitWortPlatzhalter}}
  \end{center}}

% --- Links und Querverweise (immer zuletzt) --------------------------------
\usepackage[
  hidelinks,
  bookmarks=true,
  bookmarksnumbered=true,
  pdfstartview=FitH,
  pdftitle={\kitTitel},
  pdfauthor={\kitAutor}
]{hyperref}
\ifkitEnglisch
  \usepackage[english,capitalise,nameinlink]{cleveref}
\else
  \usepackage[ngerman,capitalise,nameinlink]{cleveref}
  \crefname{equation}{Gleichung}{Gleichungen}
  \Crefname{equation}{Gleichung}{Gleichungen}
\fi

% --- Eigene Ergänzungen der Nutzerin/des Nutzers ---------------------------
\IfFileExists{eigene-praeambel.tex}{\input{eigene-praeambel.tex}}{}
